%!TeX root=MemoriaTFG.tex

\chapter{Tablas detalladas de resultados}
\label{anexo:tablas}
\label{cap:anexob}

Este anexo recoge desgloses y tablas adicionales a las del
capítulo~\ref{cap:resultados}: cifras por clase de los protocolos
de sondeo lineal y ajuste fino sobre HAM10000 y PAD-UFES-20, la
caracterización completa por fototipo cutáneo en la formulación de
$114$ patologías sobre Fitzpatrick17k, y la tabla extendida
de evaluación cruzada \emph{Light}~$\leftrightarrow$~\emph{Dark}.

\section{Desglose por clase sobre HAM10000}
\label{sec:anexob-ham}

La tabla~\ref{tab:anexob-perclass-gpt4o} reporta el desglose por
clase de la evaluación zero-shot con GPT-4o sobre HAM10000,
ampliada respecto a la presentada en
sección~\ref{sec:llm-results}. La columna ``Predicciones'' indica el
número total de imágenes asignadas a cada clase por el modelo,
revelando el sesgo léxico discutido en sección~\ref{sec:disc-llm}.

\begin{table}[H]
\centering
\footnotesize
\setlength{\tabcolsep}{4pt}
\renewcommand{\arraystretch}{0.95}
\caption{Desglose por clase del modelo GPT-4o zero-shot sobre
HAM10000. \emph{Recall}: proporción de imágenes de la clase
correctamente identificadas. \emph{Predicciones}: número de
imágenes asignadas a la clase por el modelo (debe contrastarse
con $N_{\mathrm{test}}$ para detectar sesgos sistemáticos).}
\label{tab:anexob-perclass-gpt4o}
\begin{tabular}{lrrrr}
\hline
Clase                       & $N_{\mathrm{test}}$ & Recall (GPT-4o) & Predicciones (GPT-4o) & Recall (GPT-4o-mini) \\
\hline
actinic keratosis           & $35$    & $0{,}000$ & $2$    & $0{,}086$ \\
basal cell carcinoma        & $44$    & $0{,}477$ & ---    & $0{,}136$ \\
seborrheic keratosis        & $107$   & $0{,}308$ & ---    & $0{,}654$ \\
dermatofibroma              & $8$     & $0{,}750$ & $284$  & $0{,}000$ \\
melanoma                    & $70$    & $0{,}686$ & $299$  & $0{,}571$ \\
melanocytic nevus           & $951$   & $0{,}505$ & ---    & $0{,}196$ \\
vascular lesion             & $17$    & $0{,}529$ & ---    & $0{,}588$ \\
\hline
\end{tabular}
\end{table}

La asimetría entre $N_{\mathrm{test}}$ y predicciones es máxima en
dermatofibroma ($N_{\mathrm{test}} = 8$, predicciones $284$:
factor $35{,}5\times$) y melanoma ($N_{\mathrm{test}} = 70$,
predicciones $299$: factor $4{,}3\times$). El comportamiento es
diagnóstico de un sesgo léxico hacia categorías con presencia
elevada en el \emph{prompt}, no hacia categorías mayoritarias en el
dataset.

\section{Desglose por clase del ajuste fino MedGemma~27B sobre HAM10000}
\label{sec:anexob-medgemma}

La tabla~\ref{tab:anexob-medgemma} desglosa por clase el rendimiento
de MedGemma~27B con LoRA sobre el conjunto de test de HAM10000.

\begin{table}[H]
\centering
\footnotesize
\setlength{\tabcolsep}{4pt}
\renewcommand{\arraystretch}{0.95}
\caption{Desglose por clase de MedGemma~27B con LoRA sobre
HAM10000 (\emph{precision}, \emph{recall}, F1).}
\label{tab:anexob-medgemma}
\begin{tabular}{lrrrr}
\hline
Clase                       & $N$ test & Precision & Recall    & F1 \\
\hline
actinic keratosis           & $35$     & $0{,}32$  & $0{,}26$  & $0{,}29$ \\
basal cell carcinoma        & $44$     & $0{,}34$  & $0{,}25$  & $0{,}29$ \\
dermatofibroma              & $8$      & $0{,}00$  & $0{,}00$  & $0{,}00$ \\
melanocytic nevus           & $951$    & $0{,}85$  & $0{,}99$  & $0{,}92$ \\
melanoma                    & $70$     & $0{,}38$  & $0{,}36$  & $0{,}37$ \\
seborrheic keratosis        & $107$    & $0{,}67$  & $0{,}04$  & $0{,}07$ \\
vascular lesion             & $17$     & $0{,}00$  & $0{,}00$  & $0{,}00$ \\
\hline
\end{tabular}
\end{table}

Tras la adaptación con LoRA, el modelo concentra el rendimiento en
la clase mayoritaria (nevus melanocítico: \emph{recall}
$0{,}99$) a costa de las clases menos representadas
(dermatofibroma y vascular lesion: F1 $0{,}00$). Esta
concentración es coherente con la accuracy de $0{,}802$
reportada en la tabla~\ref{tab:llm-comparativa} pese a una BAcc
de sólo $0{,}270$.

\section{Equidad sobre Fitzpatrick17k: formulación de 114 patologías}
\label{sec:anexob-fst-114}

\begin{figure}[H]
\centering
\includegraphics[width=0.78\textwidth]{fig_fairness_fst.png}
\caption{Rendimiento por fototipo Fitzpatrick sobre la
formulación de $114$ patologías de Fitzpatrick17k con PanDerm
Large + sondeo lineal. El rendimiento crece de FST~I a FST~V y
desciende bruscamente en FST~VI.}
\label{fig:anexob-fairness-fst}
\end{figure}

La figura~\ref{fig:anexob-fairness-fst} y la
tabla~\ref{tab:anexob-fst114} reportan el rendimiento por
fototipo Fitzpatrick de PanDerm Large sobre la formulación de
$114$ patologías (grano fino del dataset), complementaria a la
formulación de tres clases de malignidad presentada en la
tabla~\ref{tab:equidad-fototipo}.

\begin{table}[H]
\centering
\footnotesize
\setlength{\tabcolsep}{4pt}
\renewcommand{\arraystretch}{0.95}
\caption{Rendimiento por subgrupo Fitzpatrick sobre la
formulación de $114$ patologías de Fitzpatrick17k (PanDerm
Large + sondeo lineal).}
\label{tab:anexob-fst114}
\begin{tabular}{lrrrrr}
\hline
Fototipo (FST) & $N$ test & Acc       & BAcc      & W-F1      & AUROC \\
\hline
I              & $299$    & $0{,}515$ & $0{,}522$ & $0{,}509$ & --- \\
II             & $470$    & $0{,}519$ & $0{,}524$ & $0{,}513$ & --- \\
III            & $325$    & $0{,}640$ & $0{,}590$ & $0{,}636$ & --- \\
IV             & $261$    & $0{,}670$ & $0{,}602$ & $0{,}654$ & --- \\
V              & $169$    & $0{,}692$ & $0{,}614$ & $0{,}680$ & --- \\
VI             & $73$     & $0{,}479$ & $0{,}474$ & $0{,}462$ & --- \\
\hline
\multicolumn{2}{l}{\emph{Global}} & $0{,}581$ & $0{,}549$ & $0{,}576$ & $0{,}968$ \\
\hline
\end{tabular}
\end{table}

La distribución es notablemente distinta a la formulación de tres
clases de malignidad: el rendimiento crece desde FST~I a FST~V y
cae bruscamente en FST~VI. El \emph{gap} entre FST~V (mejor) y
FST~VI (peor) es de $0{,}213$\,pp accuracy. La interpretación
sugiere que el grano fino de $114$ patologías expone una capa
de información clínica donde la representación visual de las
lesiones sobre fototipos intermedios y altos
(III, IV, V) supera a la de fototipos extremos
(I, II y VI), mientras que la formulación gruesa de tres
clases de malignidad refleja el patrón general
peor-en-fototipos-elevados.

\section{Evaluación cruzada Light/Dark sobre 114 patologías}
\label{sec:anexob-cross-114}

La tabla~\ref{tab:anexob-cross114} reporta el experimento cruzado
de entrenamiento/evaluación entre subgrupos \emph{Light} (FST~I-III)
y \emph{Dark} (FST~IV-VI) sobre la formulación de $114$
patologías de Fitzpatrick17k con PanDerm Large. \emph{All} agrupa
el conjunto de entrenamiento completo. Las celdas de número de
clases en evaluación reflejan que algunas patologías no aparecen
en la partición \emph{Dark}, lo que reduce el número efectivo de
clases sobre las que se evalúa.

\begin{table}[H]
\centering
\footnotesize
\setlength{\tabcolsep}{4pt}
\renewcommand{\arraystretch}{0.95}
\caption{Evaluación cruzada \emph{Light}~$\leftrightarrow$~\emph{Dark}
sobre $114$ patologías de Fitzpatrick17k (PanDerm Large + sondeo
lineal). $C$: número de clases efectivas en la evaluación.}
\label{tab:anexob-cross114}
\begin{tabular}{lrrrrrr}
\hline
Configuración   & $N$ train & $N$ test & $C$ & Acc       & BAcc      & AUROC \\
\hline
Light $\to$ Light & $8\,845$ & $1\,094$ & $114$ & $0{,}541$ & $0{,}528$ & $0{,}962$ \\
Light $\to$ Dark  & $8\,028$ & $503$    & $100$ & $0{,}473$ & $0{,}453$ & $0{,}944$ \\
Dark $\to$ Dark   & $3\,807$ & $503$    & $100$ & $0{,}616$ & $0{,}519$ & $0{,}948$ \\
Dark $\to$ Light  & $3\,958$ & $1\,094$ & $114$ & $0{,}351$ & $0{,}325$ & $0{,}906$ \\
All $\to$ Light   & $12\,803$& $1\,094$ & $114$ & $0{,}559$ & $0{,}547$ & $0{,}965$ \\
All $\to$ Dark    & $11\,835$& $503$    & $100$ & $0{,}658$ & $0{,}614$ & $0{,}970$ \\
\hline
\end{tabular}
\end{table}

La asimetría más notable se observa en \emph{Dark}~$\to$~\emph{Light}
(BAcc $0{,}325$, AUROC $0{,}906$), configuración con menor
volumen de entrenamiento ($3\,958$ imágenes) sobre un conjunto
de evaluación mayor ($1\,094$ imágenes) y con $114$ clases
efectivas. La accuracy desciende desde $0{,}616$ en
\emph{Dark}~$\to$~\emph{Dark} a $0{,}351$ en
\emph{Dark}~$\to$~\emph{Light}, lo que cuantifica la pérdida
sistemática asociada al cambio de distribución de fototipo a la
inversa de la dirección habitual reportada en la literatura.

\section{Resumen comparativo por paradigma sobre HAM10000}
\label{sec:anexob-paradigmas}

La tabla~\ref{tab:anexob-paradigmas} amplía la
tabla~\ref{tab:llm-comparativa} del cuerpo con los costes de
entrenamiento e inferencia agregados por paradigma sobre HAM10000.

\begin{table}[H]
\centering
\footnotesize
\setlength{\tabcolsep}{4pt}
\renewcommand{\arraystretch}{0.95}
\caption{Comparativa por paradigma sobre HAM10000. Datos
etiquetados: tamaño del conjunto de entrenamiento empleado. Coste
indicativo en EUR para los modelos de acceso comercial (cifras
$2026$); $0$\,EUR para los modelos locales.}
\label{tab:anexob-paradigmas}
\begin{tabular}{lllrrr}
\hline
Paradigma                            & Modelo                & Datos etiq.       & Acc       & AUROC     & Coste eval.\ \\
\hline
ZS (\emph{prompts} genéricos)        & DermLIP v1/v2         & $0$               & $0{,}086$ & $0{,}366$ & $0$ \\
ZS (\emph{prompts} A3 Derm1M)        & DermLIP original      & $0$               & $0{,}427$ & $0{,}854$ & $0$ \\
ZS LLM multimodal                    & GPT-4o                & $0$               & $0{,}485$ & ---       & $\sim$\$2{,}71 \\
ZS LLM multimodal                    & MedGemma~27B          & $0$               & $0{,}665$ & ---       & $0$ \\
LP ($1\%$ entrenam.)                 & PanDerm Large         & $82$              & $0{,}850$ & $0{,}918$ & $0$ \\
LP ($100\%$ entrenam.)               & PanDerm Base          & $8\,207$          & $0{,}853$ & $0{,}892$ & $0$ \\
LP ($100\%$ entrenam.)               & PanDerm Large         & $8\,207$          & $0{,}888$ & $0{,}954$ & $0$ \\
LP ($100\%$ entrenam.)               & SigLIP-Large          & $8\,207$          & $0{,}900$ & $0{,}971$ & $0$ \\
FT LoRA LLM                          & MedGemma~27B + LoRA   & $8\,207$          & $0{,}802$ & ---       & $0$ \\
FT (receta paper) + TTA              & PanDerm Base          & $8\,207$          & $0{,}920$ & $0{,}978$ & $0$ \\
\hline
\end{tabular}
\end{table}

Las figuras~\ref{fig:anexob-tradeoff}, \ref{fig:anexob-l1-dist} y
\ref{fig:anexob-l2-acc} caracterizan el dataset armonizado de
$72\,654$ imágenes empleado en el clasificador unificado: el
compromiso entre número de clases y \emph{accuracy}, la distribución
de imágenes por categoría L1 y la \emph{accuracy} por clase L2,
respectivamente.

\begin{figure}[H]
\centering
\includegraphics[width=0.85\textwidth]{tradeoff_classes_accuracy.png}
\caption{Compromiso entre número de clases y accuracy sobre el
dataset armonizado. La accuracy decrece conforme aumenta el
número de categorías; el nivel L2 ($43$ clases) constituye el
punto operativo más equilibrado entre granularidad clínica y
rendimiento agregado.}
\label{fig:anexob-tradeoff}
\end{figure}

\begin{figure}[H]
\centering
\includegraphics[width=0.85\textwidth]{unified_l1_distribution.png}
\caption{Distribución de imágenes por categoría L1 sobre el
dataset armonizado de $72\,654$ imágenes. Predominio de
patología tumoral y patología inflamatoria; patología infecciosa
y genodermatosis representan la cola larga.}
\label{fig:anexob-l1-dist}
\end{figure}

\begin{figure}[H]
\centering
\includegraphics[width=0.95\textwidth]{unified_l2_perclass_accuracy.png}
\caption{Accuracy por clase L2 ($43$ entradas) sobre el
clasificador unificado del dataset armonizado. Las clases con
mayor accuracy corresponden a categorías con más
representación en el conjunto de entrenamiento, mientras que la
cola larga presenta variabilidad asociada al tamaño muestral.}
\label{fig:anexob-l2-acc}
\end{figure}

\section{Composición del conjunto de entrenamiento del SAE Large}
\label{sec:anexob-sae-train}

El SAE Large utilizado en el módulo M3 del prototipo y descrito en
detalle en el Anexo~\ref{cap:anexof} se entrena sobre la
concatenación de siete datasets dermatológicos. La
tabla~\ref{tab:anexob-sae-train} reporta la composición exacta.

\begin{table}[H]
\centering
\footnotesize
\setlength{\tabcolsep}{4pt}
\caption{Composición del conjunto de entrenamiento del Sparse
Autoencoder Large ($1\,024 \to 16\,384$ \emph{features}).}
\label{tab:anexob-sae-train}
\begin{tabular}{lr}
\hline
Dataset            & Imágenes \\
\hline
HAM10000           & $10\,015$ \\
BCN20000           & $12\,413$ \\
Dermnet            & $19\,559$ \\
PAD-UFES-20        & $2\,298$ \\
Fitzpatrick17k     & $3\,887$ \\
DDI                & $647$ \\
HIBA               & $1\,635$ \\
\hline
\textbf{Total}     & $\mathbf{50\,454}$ \\
\hline
\end{tabular}
\end{table}

\section{Segmentación: generalización fuera de dominio (Dice e IoU)}
\label{sec:anexob-seg-ood}

La tabla~\ref{tab:seg-detalle} amplía el resultado de segmentación de
la sección~\ref{sec:seg-results} con la desviación estándar y el IoU
sobre los dos datasets de transferencia (ISIC2017 y PH2).

\begin{table}[H]
\centering
\footnotesize
\setlength{\tabcolsep}{4pt}
\caption{Generalización fuera de dominio: Dice e IoU (media $\pm$
desv.~estándar) sobre ISIC2017 ($N = 600$) y PH2 ($N = 200$).}
\label{tab:seg-detalle}
\begin{tabular}{llrr}
\hline
Modelo                              & Dataset      & Dice              & IoU \\
\hline
SAM2 \emph{zero-shot} (bbox)        & ISIC2017     & $0{,}797 \pm 0{,}158$ & $0{,}686 \pm 0{,}183$ \\
SAM2 \emph{zero-shot} (bbox)        & PH2          & $0{,}886 \pm 0{,}199$ & $0{,}830 \pm 0{,}198$ \\
SAM2.1-L (decoder FT) ISIC2018      & ISIC2017     & $0{,}945 \pm 0{,}046$ & $0{,}898 \pm 0{,}065$ \\
SAM2.1-L (decoder FT) ISIC2018      & PH2          & $0{,}960 \pm 0{,}022$ & $0{,}925 \pm 0{,}039$ \\
CAE-seg (PanDerm-L) FT ISIC2018     & ISIC2017     & $0{,}879 \pm 0{,}139$ & $0{,}803 \pm 0{,}166$ \\
CAE-seg (PanDerm-L) FT ISIC2018     & PH2          & $0{,}924 \pm 0{,}066$ & $0{,}865 \pm 0{,}098$ \\
\hline
\end{tabular}
\end{table}
